% !PRINTED 691
\weberpage{691}
%<<CONT
and $2\nu - n$ real ones occur. Accordingly, of the $n$ functions $y$ [§. 189, (2)]

\weq{2}{y_k = \alpha_{1,k}x_1 + \alpha_{2,k}x_2 + \cdots + \alpha_{n,k}x_n.}

$2\nu - n$ are real and $2n-2\nu$ conjugate in pairs. By the $n$ variables $u$ we shall understand real linear combinations of the $x$, namely the $2\nu-n$ real $y$ and the $2n-2\nu$ parts of the conjugate imaginary pairs of the $y$, so that, when $y, y'$ is an imaginary pair, one sets
\[
2u_1 = y+y', \quad 2iu_2 = y-y'
\]
or

\weq{3}{y = u_1+iu_2, \quad y' = u_1-iu_2.}

If, in the functional determinant of the $y$, one applies to each pair of conjugate imaginary rows twice the elementary determinant theorem on the addition of rows (Bd. I, §. 22), there results

\[
\begin{aligned}
\Sigma \pm \frac{\partial u_1}{\partial x_1}\frac{\partial u_2}{\partial x_2}\cdots\frac{\partial u_n}{\partial x_n}
&= \frac{1}{(-2i)^{n-\nu}}\Sigma\pm\frac{\partial y_1}{\partial x_1}\frac{\partial y_2}{\partial x_2}\cdots\frac{\partial y_n}{\partial x_n} \\
&= \frac{A}{(-2i)^{n-\nu}},
\end{aligned}
\]

in which $A$ has the meaning fixed in §. 189, (3).

The field discriminant is positive or negative according as the number of conjugate imaginary pairs, that is $n-\nu$, is even or odd; and accordingly, if the absolute value of the field discriminant is denoted by $\pm\varDelta$, one obtains from §. 189, (4), (16):

\weq{4}{\mathrm{Lim}\,\frac{T}{t} = \frac{2^{n-\nu}}{w\,N(\mathfrak{a})\sqrt{\pm\varDelta}}\int\int\cdots\int du_1\,du_2\ldots du_n.}

Thus the integral

\[
U = \int\int\cdots\int du_1\,du_2\ldots du_n
\]

still remains to be treated.

From the equations §. 189, (10) one obtains, for every system of values of the variables $\xi_1, \xi_2, \ldots, \xi_\nu$ that satisfies the conditions (15), a definite finite system of values for $z_1, z_2, \ldots, z_\nu$, and from these, by §. 189, (7), an absolute value different from zero for every $y$. If, in the case of an imaginary pair (3), one sets

\weq{5}{u_1 = e^{1/2\,z}\cos\vartheta, \quad u_2 = e^{1/2\,z}\sin\vartheta,}

\weq{6}{y_1y_1' = u_1^2+u_2^2 = e^z,}
%CONT>>
